% GENERATED FILE. Edit canonical lessons, then run npm run generate:books.
% canonical-source-hash: fnv1a64:795c7add606a6220
% canonical-lessons: AR-C297-idhaa, AR-C297-qanah, AR-C297-midhya, AR-C297-adhaa, AR-C297-ala-al-hawa

\chapter{On the radio}
\label{ch:ar-on-the-radio}
\hlchaptermodality{297}

\begin{chapteropening}
\textbf{By the end of this chapter:} \emph{I can say a radio station, a channel, a radio set, to broadcast and on air.}

Complete the last lesson of chapter 297: I can say a radio station, a channel, a radio set, to broadcast and on air.
\end{chapteropening}

\section[\texorpdfstring{\ar{إذاعة} (idhāʿa)}{إذاعة (idhāʿa)}]{\ar{إذاعة} (idhāʿa) — a radio station; broadcasting}

\label{lesson:AR-C297-idhaa}



\begin{quote}
Before the new one: say the Arabic for to edit, then the Arabic for a bulletin.
\end{quote}

\subsection*{You'll want to know: \ar{إذاعة}}
\textbf{\ar{إذاعة}} — \emph{idhāʿa} — \textquotedblleft{}a radio station; broadcasting\textquotedblright{}.

\begin{quote}
\ar{تعمل} \ar{خالتي} \ar{في} \ar{إذاعة} \ar{صغيرة}.
\end{quote}

\begin{quote}
\emph{taʿmalu khālatī fī idhāʿa ṣaghīra.}
\end{quote}

\begin{quote}
My aunt works at a small radio station.
\end{quote}

\begin{cousinweb}
\textbf{\ar{إذاعة}} (\emph{idhāʿa}) is usually traced to the native Arabic root \emph{dh-y-ʿ}, said of news that spreads. \emph{Idhāʿa} is the verbal noun of Form IV \emph{adhāʿa}, \textquotedblleft{}to make known\textquotedblright{}, and now names radio.
\end{cousinweb}

\begin{grammarlens}[title={What you've built}]
One more word for telling what happened.
\end{grammarlens}

\subsection*{Your turn}
\begin{itemize}
\raggedright
  \item \emph{Say it:} \emph{idhāʿa}
  \item \emph{Say it:} \emph{idhāʿa}, once more
  \item \emph{Say it:} \emph{taʿmalu khālatī fī idhāʿa ṣaghīra.} — My aunt works at a small radio station.
  \item \emph{Recall:} say the Arabic for to edit, then the Arabic for a bulletin, then say \emph{idhāʿa} again
\end{itemize}

\begin{tcolorbox}[breakable,colback=teal!4,colframe=teal!35!black,title={Before you move on}]
What does \ar{إذاعة} mean? (\textquotedblleft{}A radio station; broadcasting\textquotedblright{}.) Say it once more.
\end{tcolorbox}
\section[\texorpdfstring{\ar{قناة} (qanāh)}{قناة (qanāh)}]{\ar{قناة} (qanāh) — a channel; a canal}

\label{lesson:AR-C297-qanah}



\begin{quote}
Before the new one: say the Arabic for a bulletin, then the Arabic for a radio station.
\end{quote}

\subsection*{You'll want to know: \ar{قناة}}
\textbf{\ar{قناة}} — \emph{qanāh} — \textquotedblleft{}a channel; a canal\textquotedblright{}.

\begin{quote}
\ar{تبث} \ar{هذه} \ar{القناة} \ar{الأخبار} \ar{صباحا} \ar{ومساء}.
\end{quote}

\begin{quote}
\emph{tabuththu hādhihi al-qanāh al-akhbār ṣabāḥan wa-masāʾan.}
\end{quote}

\begin{quote}
This channel broadcasts the news morning and evening.
\end{quote}

Plural \ar{قنوات}. A television channel and a water canal are one word.

\begin{cousinweb}
\textbf{\ar{قناة}} (\emph{qanāh}) is built on the native Arabic root \emph{q-n-w} (Lane), where \emph{qanāh} is a straight reed or spear shaft, and then a channel dug for water. The television sense is modern.
\end{cousinweb}

\begin{grammarlens}[title={What you've built}]
One more word for telling what happened.
\end{grammarlens}

\subsection*{Your turn}
\begin{itemize}
\raggedright
  \item \emph{Say it:} \emph{qanāh}
  \item \emph{Say it:} \emph{qanāh}, once more
  \item \emph{Say it:} \emph{tabuththu hādhihi al-qanāh al-akhbār ṣabāḥan wa-masāʾan.} — This channel broadcasts the news morning and evening.
  \item \emph{Recall:} say the Arabic for a bulletin, then the Arabic for a radio station, then say \emph{qanāh} again
\end{itemize}

\begin{tcolorbox}[breakable,colback=teal!4,colframe=teal!35!black,title={Before you move on}]
What does \ar{قناة} mean? (\textquotedblleft{}A channel; a canal\textquotedblright{}.) Say it once more.
\end{tcolorbox}
\section[\texorpdfstring{\ar{مذياع} (midhyāʿ)}{مذياع (midhyāʿ)}]{\ar{مذياع} (midhyāʿ) — a radio set}

\label{lesson:AR-C297-midhya}



\begin{quote}
Before the new one: say the Arabic for a radio station, then the Arabic for a channel.
\end{quote}

\subsection*{You'll want to know: \ar{مذياع}}
\textbf{\ar{مذياع}} — \emph{midhyāʿ} — \textquotedblleft{}a radio set\textquotedblright{}.

\begin{quote}
\ar{كان} \ar{جدي} \ar{يسمع} \ar{المذياع} \ar{كل} \ar{صباح}.
\end{quote}

\begin{quote}
\emph{kāna jaddī yasmaʿu al-midhyāʿ kulla ṣabāḥ.}
\end{quote}

\begin{quote}
My grandfather used to listen to the radio every morning.
\end{quote}

Many speakers say \ar{راديو}; \ar{مذياع} is the standard word.

\begin{cousinweb}
\textbf{\ar{مذياع}} (\emph{midhyāʿ}) is usually traced to the native Arabic root \emph{dh-y-ʿ}, \textquotedblleft{}to spread, to become known\textquotedblright{}. The \emph{mifʿāl} pattern names a tool, so \emph{midhyāʿ} is \textquotedblleft{}the spreader\textquotedblright{}: a modern coinage for the radio set.
\end{cousinweb}

\begin{grammarlens}[title={What you've built}]
One more word for telling what happened.
\end{grammarlens}

\subsection*{Your turn}
\begin{itemize}
\raggedright
  \item \emph{Say it:} \emph{midhyāʿ}
  \item \emph{Say it:} \emph{midhyāʿ}, once more
  \item \emph{Say it:} \emph{kāna jaddī yasmaʿu al-midhyāʿ kulla ṣabāḥ.} — My grandfather used to listen to the radio every morning.
  \item \emph{Recall:} say the Arabic for a radio station, then the Arabic for a channel, then say \emph{midhyāʿ} again
\end{itemize}

\begin{tcolorbox}[breakable,colback=teal!4,colframe=teal!35!black,title={Before you move on}]
What does \ar{مذياع} mean? (\textquotedblleft{}A radio set\textquotedblright{}.) Say it once more.
\end{tcolorbox}
\section[\texorpdfstring{\ar{أذاع} (adhāʿa)}{أذاع (adhāʿa)}]{\ar{أذاع} (adhāʿa) — to put on air; to make known}

\label{lesson:AR-C297-adhaa}



\begin{quote}
Before the new one: say the Arabic for a channel, then the Arabic for a radio set.
\end{quote}

\subsection*{You'll want to know: \ar{أذاع}}
\textbf{\ar{أذاع}} — \emph{adhāʿa} — \textquotedblleft{}to put on air; to make known\textquotedblright{}.

\begin{quote}
\ar{أذاعت} \ar{الإذاعة} \ar{أغنية} \ar{قديمة}.
\end{quote}

\begin{quote}
\emph{adhāʿat al-idhāʿa ughniya qadīma.}
\end{quote}

\begin{quote}
The station broadcast an old song.
\end{quote}

Present: \ar{يذيع}, \emph{yudhīʿ}.

\begin{cousinweb}
\textbf{\ar{أذاع}} (\emph{adhāʿa}) is usually traced to the native Arabic root \emph{dh-y-ʿ}, said of news that spreads. It is Form IV, \textquotedblleft{}to make known\textquotedblright{}, and the verb of \ar{إذاعة}.
\end{cousinweb}

\begin{grammarlens}[title={What you've built}]
One more word for telling what happened.
\end{grammarlens}

\subsection*{Your turn}
\begin{itemize}
\raggedright
  \item \emph{Say it:} \emph{adhāʿa}
  \item \emph{Say it:} \emph{adhāʿa}, once more
  \item \emph{Say it:} \emph{adhāʿat al-idhāʿa ughniya qadīma.} — The station broadcast an old song.
  \item \emph{Recall:} say the Arabic for a channel, then the Arabic for a radio set, then say \emph{adhāʿa} again
\end{itemize}

\begin{tcolorbox}[breakable,colback=teal!4,colframe=teal!35!black,title={Before you move on}]
What does \ar{أذاع} mean? (\textquotedblleft{}To put on air; to make known\textquotedblright{}.) Say it once more.
\end{tcolorbox}
\section[\texorpdfstring{\ar{على} \ar{الهواء} (ʿalā al-hawāʾ)}{على الهواء (ʿalā al-hawāʾ)}]{\ar{على} \ar{الهواء} (ʿalā al-hawāʾ) — on air, live}

\label{lesson:AR-C297-ala-al-hawa}



\begin{quote}
Before the new one: say the Arabic for a radio set, then the Arabic for to put on air.
\end{quote}

\subsection*{You'll want to know: \ar{على} \ar{الهواء}}
\textbf{\ar{على} \ar{الهواء}} — \emph{ʿalā al-hawāʾ} — \textquotedblleft{}on air, live\textquotedblright{}.

\begin{quote}
\ar{تكلم} \ar{الوزير} \ar{على} \ar{الهواء} \ar{مساء} \ar{أمس}.
\end{quote}

\begin{quote}
\emph{takallama al-wazīr ʿalā al-hawāʾ masāʾa ams.}
\end{quote}

\begin{quote}
The minister spoke on air yesterday evening.
\end{quote}

\ar{على} \ar{الهواء} \ar{مباشرة}: live on air.

\begin{cousinweb}
\textbf{\ar{على} \ar{الهواء}} (\emph{ʿalā al-hawāʾ}) is \emph{ʿalā}, \textquotedblleft{}on\textquotedblright{}, with \emph{hawāʾ}, \textquotedblleft{}air\textquotedblright{}, which is built on the native Arabic root \emph{h-w-y} (Lane). The phrase is usually traced to the English broadcasting phrase on the air.
\end{cousinweb}

\begin{grammarlens}[title={What you've built}]
One more word for telling what happened.
\end{grammarlens}

\subsection*{Your turn}
\begin{itemize}
\raggedright
  \item \emph{Say it:} \emph{ʿalā al-hawāʾ}
  \item \emph{Say it:} \emph{ʿalā al-hawāʾ}, once more
  \item \emph{Say it:} \emph{takallama al-wazīr ʿalā al-hawāʾ masāʾa ams.} — The minister spoke on air yesterday evening.
  \item \emph{Recall:} say the Arabic for a radio set, then the Arabic for to put on air, then say \emph{ʿalā al-hawāʾ} again
\end{itemize}

\begin{tcolorbox}[breakable,colback=teal!4,colframe=teal!35!black,title={Before you move on}]
What does \ar{على} \ar{الهواء} mean? (\textquotedblleft{}On air, live\textquotedblright{}.) Say it once more.
\end{tcolorbox}
